\section{问题重述}

本题考虑由六个区域构成的异构算力系统。系统接收实时推理、批量推理和 AI 训练三类任务，每项任务均具有来源区域、到达时刻、GPU 需求、预计持续时间、最大允许网络时延及最晚完成时点等属性。各区域具有不同的 GPU 容量、IT 功率上限、设施功率上限和 PUE，同时受到逐时购售电价格、碳强度、新能源出力、电网接入边界及储能运行边界的共同影响。需要在任务不可拆分、不可抢占且不可中途迁移的条件下，确定任务执行区域、开工时刻以及能源侧运行策略。

系统采用分层时域。第 $0$ 至 $2399$ 小时为任务到达主时域，第 $2400$ 至 $2405$ 小时不再接收新任务，仅用于完成此前到达的弹性任务；第 $2406$ 小时不得执行任务，仅用于电力平衡和储能终端状态结算。实时推理任务必须在到达时刻开工，批量推理和 AI 训练任务可在允许窗口内延迟执行，但完成时点不得超过任务截止时点与系统终点中的较早者。任务持续时间以分钟记录，因此其 GPU 占用和 AI IT 功率需依据任务与各小时的实际重叠时长计算。

问题一要求统计不同区域、不同任务类型的 GPU 需求特征，构建无时间泄漏的短期预测模型，并针对第 $2376$ 至 $2399$ 小时的实际到达任务给出基础调度方案及末端结清安排。问题二要求基于完整实际任务和逐时电力参数进行碳感知调度，在满足算力、功率、时延及完成时限约束的同时，评价运行成本、碳排放、网络时延和新能源利用率。问题三固定任务负荷，仅优化储能充放电、购售电与新能源分配，并比较相应的成本、碳排放、峰值净购电和负荷波动。问题四进一步联合任务调度和能源调度，在成本、碳排放、网络时延、服务质量、新能源利用率及区域峰值净购电之间进行权衡，并考察碳约束、电价机制和新能源波动变化对系统策略的影响。
